% Provenance (never printed): NSI premiere "Langages et programmation: utilisation de bibliotheques, documentation;
% modularite; diversite et unite des langages". See sources/book2/BRIEFS.md.
\chapter{Modules, Interfaces and Libraries}\label{ch:g11:modules-and-interfaces}

Twenty programs of the town hall use the same small \omterm{def:g11:modules-and-interfaces:module}{module} of dates: the school calendar, the \omterm{def:g11:modules-and-interfaces:library}{library}'s loans, the
market's opening days. One summer the \omterm{def:g11:modules-and-interfaces:module}{module} is rewritten and becomes thousands of times faster on large calendars.
None of the twenty programs is changed, none of them breaks, and their authors only notice that things run faster.

That is possible because the twenty programs relied only on what the \omterm{def:g11:modules-and-interfaces:module}{module} promised, never on how it kept its
dates. This chapter splits programs into \omterm{def:g11:modules-and-interfaces:module}{modules}, separates the promise from the way it is kept, reads the
documentation of Python's own \omterm{def:g11:modules-and-interfaces:module}{modules}, and then changes an \omterm{def:g11:modules-and-interfaces:interface}{implementation} behind its users' backs, with the tests to
prove nothing broke.

\begin{recall}
A function has parameters and returns a value (the first book of the series). A specification states what a function
must do, and tests check it (chapter \ref{ch:g10:specifications-and-tests}); a set answers \code{x in s} without a
scan (chapter \ref{ch:g10:tuples-sets-comprehensions}); the \omterm{def:g11:cost-of-an-algorithm:halvings}{number of halvings} $h(n)$ and the cost of searching
(chapter \ref{ch:g11:cost-of-an-algorithm}).
\end{recall}

\section{Splitting a program into modules}\label{sec:g11:modules-and-interfaces:modules}

\begin{definition}[Module]\label{def:g11:modules-and-interfaces:module}
A \emph{module}\index{module} is a file of code whose names, functions, constants and others, other programs can use by
importing it. In Python every file \code{name.py} is a module named \code{name}.
\end{definition}

\omcode{code/grade-11/09-modules-and-interfaces/listings/geometry.py}{1}{23}{A \omterm{def:g11:modules-and-interfaces:module}{module} of plane geometry: two functions
and a helper.}\label{lst:g11:modules-and-interfaces:geometry}

\code{import geometry} runs the file once and makes the single name \code{geometry} available; its contents are
reached through it, as \code{geometry.distance}. \code{from geometry import perimeter} copies one name into the
program, to be used without the prefix.

\omcode{code/grade-11/09-modules-and-interfaces/listings/use_geometry.py}{1}{8}{Two ways of importing, and the names a
\omterm{def:g11:modules-and-interfaces:module}{module} holds.}\label{lst:g11:modules-and-interfaces:use-geometry}

\omoutput{code/grade-11/09-modules-and-interfaces/outputs/use_geometry.txt}{The \omterm{def:g11:modules-and-interfaces:module}{module} also holds \code{math}, which it
imported for itself, and the helper \code{_square}.}\label{out:g11:modules-and-interfaces:use-geometry}

Python looks for a \omterm{def:g11:modules-and-interfaces:module}{module} first in the directory of the program being run, then in the directories listed in the
variable \code{PYTHONPATH}, then where the installation keeps its \omterm{def:g11:modules-and-interfaces:module}{modules}; the list is in \code{sys.path}. A \omterm{def:g11:modules-and-interfaces:module}{module}
named like one of Python's own, \code{random.py} next to a program, hides the real one: a classic trap.
% ledger: g11:py-syspath

Each \omterm{def:g11:modules-and-interfaces:module}{module} has a name in the variable \code{__name__}. A file started as a program gets the name
\code{"__main__"}; imported, it gets its own name. The guard \code{if __name__ == "__main__":} keeps a demonstration
from running when the file is only imported.
% ledger: g11:py-main-name

\omcode{code/grade-11/09-modules-and-interfaces/listings/main_guard.py}{1}{10}{A \omterm{def:g11:modules-and-interfaces:module}{module} with a guarded
demonstration.}\label{lst:g11:modules-and-interfaces:main-guard}

\omoutput{code/grade-11/09-modules-and-interfaces/outputs/main_guard.txt}{Started as a program: the guard lets the
greeting through.}\label{out:g11:modules-and-interfaces:main-guard}

\omoutput{code/grade-11/09-modules-and-interfaces/outputs/import_guard.txt}{Imported by \code{import_guard.py}, which
calls \code{greet} itself: the \omterm{def:g11:modules-and-interfaces:module}{module}'s own greeting does not run.}\label{out:g11:modules-and-interfaces:import-guard}

A \omterm{def:g11:modules-and-interfaces:module}{module}'s statements run only at its first import; a second \code{import} in the same run finds it already loaded.
% ledger: g11:py-import-once

\begin{omfigure}
\begin{tikzpicture}[mod/.style={draw, rounded corners=2pt, fill=omDef!8, font=\ttfamily\footnotesize,
    minimum height=6mm, inner sep=3pt}, std/.style={mod, fill=black!6}, ->, >=stealth]
  \node[mod] (use) at (0,0) {use\_dateset.py};
  \node[mod] (tests) at (4.2,0) {dateset\_tests.py};
  \node[mod] (meas) at (9.2,0) {measure\_dateset.py};
  \node[std] (il) at (4.2,1.3) {importlib, unittest};
  \node[std] (st) at (9.2,1.3) {statistics, time};
  \node[std] (dt) at (-1.5,-1.6) {datetime};
  \node[mod] (l) at (1.4,-1.6) {dateset\_list};
  \node[mod] (s) at (4.6,-1.6) {dateset\_set};
  \node[mod] (o) at (8.0,-1.6) {dateset\_sorted};
  \node[std] (bi) at (8.0,-3.0) {bisect};
  \draw (tests) -- (il); \draw (meas) -- (st);
  \foreach \p in {use, tests, meas} { \draw (\p) -- (dt); \draw (\p) -- (l); \draw (\p) -- (s); }
  \draw (tests) -- (o); \draw (meas) -- (o);
  \draw (o) -- (bi);
\end{tikzpicture}

{\small The \omterm{def:g11:modules-and-interfaces:module}{modules} of this chapter's date programs. An arrow goes from a \omterm{def:g11:modules-and-interfaces:module}{module} to a \omterm{def:g11:modules-and-interfaces:module}{module} it imports; the grey
ones come with Python.}
\end{omfigure}

\section{The interface and the hidden implementation}\label{sec:g11:modules-and-interfaces:interface}

\begin{definition}[Interface, implementation]\label{def:g11:modules-and-interfaces:interface}
The \emph{interface}\index{interface} of a \omterm{def:g11:modules-and-interfaces:module}{module} is what its users may rely on: the names it offers and, for each, a
specification. The \emph{implementation}\index{implementation} is the code that keeps these promises, and everything
else in the \omterm{def:g11:modules-and-interfaces:module}{module}: helpers, data structures, imported \omterm{def:g11:modules-and-interfaces:module}{modules}.
\end{definition}

Python does not lock anything away: \code{geometry._square} can be called. The protection is a convention: a name
beginning with an underscore is internal, and \code{from geometry import *} does not even copy it. A name that is not
documented, or that the \omterm{def:g11:modules-and-interfaces:module}{module} only imported for itself, as \code{geometry.math}, is internal too.
% ledger: g11:pep8-names
% ledger: g11:py-import-star

The documentation of a function is its \omterm{def:g10:specifications-and-tests:docstring}{docstring}, which \code{help} prints.

\omcode{code/grade-11/09-modules-and-interfaces/listings/help_view.py}{1}{3}{Asking for the documentation of a
function.}\label{lst:g11:modules-and-interfaces:help}

\omoutput{code/grade-11/09-modules-and-interfaces/outputs/help_view.txt}{The signature and the \omterm{def:g10:specifications-and-tests:docstring}{docstring}: all a user
needs.}\label{out:g11:modules-and-interfaces:help}

\begin{omfigure}
\begin{tikzpicture}[font=\footnotesize]
  \fill[omDef!8] (0,0) rectangle (11,1.7);
  \fill[black!8] (0,-1.9) rectangle (11,0);
  \draw[very thick, omThm] (0,0) -- (11,0);
  \node[right, font=\small\bfseries] at (0.1,1.45) {interface: what users may rely on};
  \node[right, align=left] at (0.2,0.7) {\code{empty()}\quad \code{from_dates(dates)}\quad \code{add(s, d)}\quad
    \code{contains(s, d)}\quad \code{size(s)}\\ and the docstring of each};
  \node[right, font=\small\bfseries] at (0.1,-0.3) {implementation: hidden, may change};
  \node[right, align=left] at (0.2,-1.1) {\code{dateset_list}: a list scanned by \code{in}\\
    \code{dateset_set}: a Python set};
  \node[right, align=left] at (6,-1.1) {\code{dateset_sorted}: a sorted list,\\ searched by halving};
\end{tikzpicture}

{\small The \omterm{def:g11:modules-and-interfaces:interface}{interface} as a wall. Above it, five names with their specifications, the same in the three \omterm{def:g11:modules-and-interfaces:module}{modules}; below
it, three ways of keeping the dates, which users never see.}
\end{omfigure}

\section{The standard library and its documentation}\label{sec:g11:modules-and-interfaces:library}

\begin{definition}[Library, standard library]\label{def:g11:modules-and-interfaces:library}
A \emph{library}\index{library} is a collection of \omterm{def:g11:modules-and-interfaces:module}{modules} written to be used by other programs. The \emph{standard
library}\index{standard library} of a language is the library that comes with every installation of it.
\end{definition}

\begin{definition}[API]\label{def:g11:modules-and-interfaces:api}
The \emph{API}\index{API (application programming interface)}, for application programming \omterm{def:g11:modules-and-interfaces:interface}{interface}, of a \omterm{def:g11:modules-and-interfaces:library}{library}
is its \omterm{def:g11:modules-and-interfaces:interface}{interface}: the names, with their specifications, that programs may use.
\end{definition}

Python's \omterm{def:g11:modules-and-interfaces:library}{standard library} has hundreds of \omterm{def:g11:modules-and-interfaces:module}{modules}, and its documentation describes each function in a fixed shape:
the signature, what it returns, what it raises and when, examples, and in which version it changed.

\begin{omfigure}
\begin{tikzpicture}[font=\footnotesize, note/.style={font=\scriptsize, omThm, align=left}]
  \draw[fill=black!3] (0,0) rectangle (7.6,-4.3);
  \node[right, font=\ttfamily\small] at (0.15,-0.35) {statistics.median(data)};
  \node[right, text width=7.1cm] at (0.15,-1.35) {Return the median (middle value) of numeric data. When the number
    of data points is even, the median is interpolated by taking the average of the two middle values.};
  \node[right, text width=7.1cm] at (0.15,-2.3) {If \emph{data} is empty, \code{StatisticsError} is raised.};
  \node[right, font=\ttfamily] at (0.15,-3.0) {>>> median([1, 3, 5, 7])};
  \node[right, font=\ttfamily] at (0.15,-3.4) {4.0};
  \node[right, font=\itshape] at (0.15,-3.95) {Changed in version \dots};
  \foreach \y/\t in {-0.35/{signature: name and\\ parameters}, -1.35/{what it returns}, -2.3/{what it raises,\\ and when},
      -3.2/{an example to\\ run and compare}, -3.95/{the versions}}
    { \draw[omThm, ->] (9.4,\y) -- (7.7,\y); \node[note, right] at (9.45,\y) {\t}; }
\end{tikzpicture}

{\small The parts of a documentation entry, on the entry for \code{statistics.median}, shortened. Reading the
raised exceptions and the examples is part of reading the entry.}
\end{omfigure}
% ledger: g11:py-median

\omcode{code/grade-11/09-modules-and-interfaces/listings/stdlib_tour.py}{1}{11}{Four \omterm{def:g11:modules-and-interfaces:module}{modules} of the \omterm{def:g11:modules-and-interfaces:library}{standard library},
used after reading their entries.}\label{lst:g11:modules-and-interfaces:stdlib}

\omoutput{code/grade-11/09-modules-and-interfaces/outputs/stdlib_tour.txt}{A choice, two medians, the three most
common letters, and some dates.}\label{out:g11:modules-and-interfaces:stdlib}

Each line of the output is explained by an entry:
\begin{itemize}
  \item \code{random.choice} returns an element of a non-empty sequence, and raises \code{IndexError} on an empty one.
    The documentation also warns that, for a given seed, only \code{random()} itself is guaranteed to give the same
    numbers in every version of Python; \code{choice} may change.
    % ledger: g11:py-random-choice
  \item \code{statistics.median} of an even number of values is the mean of the two middle ones, here $(3 + 4) / 2$.
  \item \code{Counter.most_common} lists equal counts in the order first met: \code{b} before \code{r}.
    % ledger: g11:py-most-common
  \item \code{date.weekday()} numbers Monday $0$ to Sunday $6$: $4$ is a Friday. Subtracting two dates gives a
    \code{timedelta}, here $42$ days; years run from $1$ to $9999$.
    % ledger: g11:py-date
\end{itemize}

\begin{remark}
The documentation is written for one version of Python. The entries of this book were read for version $3.12$; an
entry that says ``changed in version'' or ``new in version'' tells what an older or newer Python does differently.
\end{remark}

\section{Changing an implementation without touching its users}\label{sec:g11:modules-and-interfaces:change}

The \omterm{def:g11:modules-and-interfaces:module}{module} of dates of the hook offers five functions. Its first \omterm{def:g11:modules-and-interfaces:interface}{implementation} keeps the dates in a list.

\omcode{code/grade-11/09-modules-and-interfaces/listings/dateset_list.py}{1}{30}{A set of dates kept in a
list.}\label{lst:g11:modules-and-interfaces:dateset-list}

The second keeps them in a Python set; only the bodies change, and they become shorter.

\omcode{code/grade-11/09-modules-and-interfaces/listings/dateset_set.py}{7}{29}{The same \omterm{def:g11:modules-and-interfaces:interface}{interface}, kept in a set
(lines 1--4, the \omterm{def:g11:modules-and-interfaces:module}{module}'s \omterm{def:g10:specifications-and-tests:docstring}{docstring}, are not shown).}\label{lst:g11:modules-and-interfaces:dateset-set}

\omcode{code/grade-11/09-modules-and-interfaces/listings/use_dateset.py}{1}{12}{A user of the \omterm{def:g11:modules-and-interfaces:interface}{interface}, run with each
\omterm{def:g11:modules-and-interfaces:interface}{implementation}.}\label{lst:g11:modules-and-interfaces:use-dateset}

\omoutput{code/grade-11/09-modules-and-interfaces/outputs/use_dateset.txt}{The same answers from both: the second
1 May was not added twice.}\label{out:g11:modules-and-interfaces:use-dateset}

The tests are written against the \omterm{def:g11:modules-and-interfaces:interface}{interface} only, and the same file tests every \omterm{def:g11:modules-and-interfaces:interface}{implementation}, named on the command
line.

\omcode{code/grade-11/09-modules-and-interfaces/listings/dateset_tests.py}{1}{33}{The tests of the \omterm{def:g11:modules-and-interfaces:interface}{interface}: one file
for every \omterm{def:g11:modules-and-interfaces:interface}{implementation}.}\label{lst:g11:modules-and-interfaces:tests}

\omoutput{code/grade-11/09-modules-and-interfaces/outputs/run_dateset_tests.txt}{The same three tests pass on both
\omterm{def:g11:modules-and-interfaces:interface}{implementations} (the line with the running time is filtered out).}\label{out:g11:modules-and-interfaces:tests}

\begin{method}[Changing an implementation safely]
\begin{enumerate}
  \item Write the \omterm{def:g11:modules-and-interfaces:interface}{interface} down: the names and the specification of each, in \omterm{def:g10:specifications-and-tests:docstring}{docstrings}.
  \item Write tests that use the \omterm{def:g11:modules-and-interfaces:interface}{interface} only, and check that they pass on the old \omterm{def:g11:modules-and-interfaces:interface}{implementation}.
  \item Write the new \omterm{def:g11:modules-and-interfaces:interface}{implementation} with the same names and specifications.
  \item Run the same tests on it. Users need change nothing but, at most, the name of the \omterm{def:g11:modules-and-interfaces:module}{module} they import.
\end{enumerate}
\end{method}

\section{Libraries outside the standard library, and what languages share}\label{sec:g11:modules-and-interfaces:outside}

\begin{inpractice}[Libraries from elsewhere]
Thousands of libraries are published outside the \omterm{def:g11:modules-and-interfaces:library}{standard library}. Python's installer, \code{pip}, fetches them from a
public index; other languages have their own installers. Before relying on one, a programmer checks its version (an
update may change the \omterm{def:g11:modules-and-interfaces:interface}{interface}), its licence (what it allows), and whether it is still maintained. This book installs
nothing: every program in it uses the \omterm{def:g11:modules-and-interfaces:library}{standard library} only.
\end{inpractice}

Languages differ in their spelling, much less in their constructions. The same few constructions appear in the toy
machine's \omterm{def:g11:von-neumann-machine:assembly}{assembly language} of chapter \ref{ch:g11:von-neumann-machine} and in Python:
\begin{center}
\begin{tabular}{lll}
  construction & toy assembly & Python \\ \hline
  assignment & \code{LOADI R0, 3} & \code{r0 = 3} \\
  sequence & one instruction after another & one line after another \\
  conditional & \code{CMP} then \code{JC} to a label & \code{if a < b:} \\
  loop & \code{JNZ loop} back to a label & \code{while r0 != 0:} \\
  call & \code{CALL square}, \code{RET} & \code{square(x)}, \code{return}
\end{tabular}
\end{center}
A compiler or an interpreter translates the right-hand column into something like the middle one.

\section{Exercises}\label{sec:g11:modules-and-interfaces:exercises}

\begin{exercise}[$\star$]\label{exo:g11:modules-and-interfaces:1}
After \code{import geometry} alone, which of these work: \code{geometry.distance((0, 0), (1, 1))},
\code{distance((0, 0), (1, 1))}, \code{geometry.math.sqrt(4)}? Which should a user avoid, even if it works?
\end{exercise}

\begin{exercise}[$\star$]\label{exo:g11:modules-and-interfaces:2}
The file \code{report.py} holds the line \code{print(__name__)}. The file \code{main.py} holds \code{import report}
then \code{print(__name__)}. What does \code{python main.py} print? And \code{python report.py}?
\end{exercise}

\begin{exercise}[$\star$]\label{exo:g11:modules-and-interfaces:3}
A \omterm{def:g11:modules-and-interfaces:module}{module} defines \code{MAX_SIZE}, \code{_cache}, \code{load}, \code{_parse} and \code{save}, and imports \code{json}.
Which names belong to its \omterm{def:g11:modules-and-interfaces:interface}{interface}?
\end{exercise}

\begin{exercise}[$\star$]\label{exo:g11:modules-and-interfaces:4}
From the entries quoted in the chapter, give without running them \code{statistics.median([7, 1, 3, 9])} and
\code{Counter("banana").most_common(2)}.
\end{exercise}

\begin{exercise}[$\star$]\label{exo:g11:modules-and-interfaces:5}
Write a \omterm{def:g11:modules-and-interfaces:module}{module} \code{temperature} with two functions, \code{to_fahrenheit(c)} and \code{to_celsius(f)}, each with a
\omterm{def:g10:specifications-and-tests:docstring}{docstring}. (Fahrenheit is Celsius times $9/5$, plus $32$.)
\end{exercise}

\begin{exercise}[$\star\star$]\label{exo:g11:modules-and-interfaces:6}
\code{stops.py} starts with \code{import lines}, defines \code{FIRST_STOP = "Market Square"}, then prints
\code{lines.LINES[0]}; \code{lines.py} starts with \code{import stops} and defines \code{LINES} from
\code{stops.FIRST_STOP}. \code{python stops.py} stops with
\texttt{AttributeError: partially initialized module \textquotesingle{}lines\textquotesingle{} has no attribute \textquotesingle{}LINES\textquotesingle{}}. Explain, and correct.
\end{exercise}

\begin{exercise}[$\star\star$]\label{exo:g11:modules-and-interfaces:7}
A program does \code{holidays = dateset_list.from_dates(days)}, then \code{holidays.sort()} and prints
\code{holidays[0]}, the earliest holiday. It breaks when the town switches to \code{dateset_set}. Whose fault is it,
and how should the earliest date be obtained?
\end{exercise}

\begin{exercise}[$\star\star$]\label{exo:g11:modules-and-interfaces:8}
Rewrite \code{geometry.perimeter} without indices, with \code{zip}, so that the same tests still pass.
\end{exercise}

\begin{exercise}[$\star\star$]\label{exo:g11:modules-and-interfaces:9}
For $n$ dates, how many dates do \code{add} and \code{contains} look at, at most, in \code{dateset_list}? What about
\code{dateset_set}?
\end{exercise}

\begin{exercise}[$\star\star$]\label{exo:g11:modules-and-interfaces:10}
What does this print, and why? \code{print(pow(2, 10))}, then \code{from math import *}, then again
\code{print(pow(2, 10))}.
\end{exercise}

\begin{exercise}[$\star\star$]\label{exo:g11:modules-and-interfaces:11}
After reading the entry for \code{Counter}, write \code{word_count(text, n)}, which returns the \code{n} most common
words of a text in lower case, with their counts, punctuation removed.
\end{exercise}

\begin{exercise}[$\star\star\star$]\label{exo:g11:modules-and-interfaces:12}
Design, in prose, the \omterm{def:g11:modules-and-interfaces:interface}{interface} of a \omterm{def:g11:modules-and-interfaces:module}{module} of queues, the waiting line where the first in is the first out: the names
of the functions and the specification of each, including what happens on an empty queue.
\end{exercise}

\begin{exercise}[$\star\star\star$]\label{exo:g11:modules-and-interfaces:13}
Write a \omterm{def:g11:modules-and-interfaces:module}{module} \code{dates_chosen} that offers the five names of the \omterm{def:g11:modules-and-interfaces:interface}{interface}, taken from \code{dateset_list} when the
environment variable \code{DATESET} is \code{list}, and from \code{dateset_set} otherwise. How can a single test file
check both?
\end{exercise}

\begin{exercise}[$\star\star\star$]\label{exo:g11:modules-and-interfaces:14}
Twenty programs call \code{contains} in about ten places each. Compare how many lines of user code change when the
\omterm{def:g11:modules-and-interfaces:interface}{implementation} changes from list to set, and when the \omterm{def:g11:modules-and-interfaces:interface}{interface} renames \code{contains} to \code{has}.
\end{exercise}

\begin{exercise}[$\star\star\star$]\label{exo:g11:modules-and-interfaces:15}
The \omterm{def:g11:modules-and-interfaces:module}{module} \code{primes} computes, at its top level, all the primes below a million, and prints
\code{computing the primes below 1000000...}. Every program that imports it, even to read a constant, prints the
message and waits. Correct the \omterm{def:g11:modules-and-interfaces:module}{module} so that the work is done only when a prime is first asked for.
\end{exercise}

\section{Problem: The Calendar Module}\label{sec:g11:modules-and-interfaces:problem}

\begin{problem}[{Weekend problem --- a module three times rewritten}]\label{pb:g11:modules-and-interfaces:1}
The town's date \omterm{def:g11:modules-and-interfaces:module}{module} now has three \omterm{def:g11:modules-and-interfaces:interface}{implementations} behind one \omterm{def:g11:modules-and-interfaces:interface}{interface}. The third keeps the dates in a sorted list
and searches it by halving. The program below measures \code{contains} in all three, for $n = 10^3$, $10^4$ and
$10^5$ dates.

\omcode{code/grade-11/09-modules-and-interfaces/listings/dateset_sorted.py}{5}{32}{The third \omterm{def:g11:modules-and-interfaces:interface}{implementation}: a sorted
list.}\label{lst:g11:modules-and-interfaces:dateset-sorted}

\omcode{code/grade-11/09-modules-and-interfaces/listings/measure_dateset.py}{1}{35}{One lookup, three
\omterm{def:g11:modules-and-interfaces:interface}{implementations}.}\label{lst:g11:modules-and-interfaces:measure}

\omoutput{code/grade-11/09-modules-and-interfaces/outputs/measure_dateset.txt}{Measured on the author's laptop, in
microseconds per lookup.}\label{out:g11:modules-and-interfaces:measure}

\begin{omfigure}
\begin{tikzpicture}
  \begin{axis}[width=10cm, height=6.5cm, xmode=log, ymode=log, xlabel={$n$ dates}, ylabel={microseconds per lookup},
      legend pos=north west, legend cell align=left, font=\footnotesize]
    \addplot[omThm, thick, mark=*] table[x=n, y=list]
      {code/grade-11/09-modules-and-interfaces/figures/measured_contains.dat};
    \addplot[omMeth, thick, mark=square*] table[x=n, y=sorted]
      {code/grade-11/09-modules-and-interfaces/figures/measured_contains.dat};
    \addplot[omDef, thick, mark=triangle*] table[x=n, y=set]
      {code/grade-11/09-modules-and-interfaces/figures/measured_contains.dat};
    \legend{list, sorted list, set}
  \end{axis}
\end{tikzpicture}

{\small \code{contains} measured on the author's laptop (\code{figures/measured_contains.py}, a separate run). On both
axes each step multiplies by $10$. The list climbs in step with $n$; the two others stay flat.}
\end{omfigure}

\medskip\noindent\textbf{Part I --- A module of dates.}
\begin{enumerate}
  \item With the leap-year rule of chapter \ref{ch:g10:specifications-and-tests}, how many days have the years $2028$,
    $2100$ and $2000$?
  \item What does \code{size} return for a set of dates holding every day of $2028$?
  \item \code{date(2027, 2, 29)} raises \code{ValueError: day is out of range for month}. Why is this a good
    behaviour for a \omterm{def:g11:modules-and-interfaces:module}{module} of dates?
  \item From the output of \code{stdlib_tour.py}, on which day of the week is 20 November 2026, and which date comes
    $100$ days later?
\end{enumerate}

\medskip\noindent\textbf{Part II --- The interface and its tests.}
\begin{enumerate}[resume]
  \item List the five names of the \omterm{def:g11:modules-and-interfaces:interface}{interface}, with the specification of each in one sentence.
  \item Which test checks that adding the same date twice keeps one copy?
  \item Why do the tests call only the five names of the \omterm{def:g11:modules-and-interfaces:interface}{interface}?
  \item What should \code{python dateset_tests.py sorted} print?
  \item Write one more test: \code{from_dates([])} gives a set of size $0$.
\end{enumerate}

\medskip\noindent\textbf{Part III --- Three implementations.}
\begin{enumerate}[resume]
  \item In \code{dateset_sorted}, which function relies on the list being sorted, and which keeps it sorted?
  \item For $n = 10^5$, how many dates does \code{contains} look at, at most, in the list, and in the sorted list?
  \item Why does \code{from_dates} in \code{dateset_list} use \code{dict.fromkeys} rather than \code{list(dates)}?
  \item Why is \code{add} slower than \code{contains} in the sorted \omterm{def:g11:modules-and-interfaces:interface}{implementation}?
\end{enumerate}

\medskip\noindent\textbf{Part IV --- Measured cost.}
\begin{enumerate}[resume]
  \item By how much is the list's time multiplied from $n = 10^3$ to $n = 10^5$? Is it what the count predicts?
  \item How do the sorted list and the set behave as $n$ grows?
  \item Predict the list's time per lookup for $n = 10^6$.
  \item A program makes a thousand lookups in $10^5$ dates. How long does it take with each \omterm{def:g11:modules-and-interfaces:interface}{implementation}?
  \item The measurement uses \code{from_dates} to build the sets. Why not $n$ calls to \code{add}?
  \item Which lines of \code{use_dateset.py} would change to use \code{dateset_sorted}?
  \item State the result: the speed-up of the set over the list at $n = 10^5$, and what happened to the tests.
\end{enumerate}
\end{problem}
